% ============================================================================
% CAPÍTULO 6: SISTEMAS GÊNICOS
% Bioinformática para Biologia de Sistemas
% ============================================================================

% ============================================================================
% TEMPLATE BASE PARA APRESENTAÇÕES EM BEAMER
% Bioinformática para Biologia de Sistemas
% ============================================================================

\documentclass[12pt, aspectratio=169]{beamer}

% ============================================================================
% PACOTES ESSENCIAIS
% ============================================================================
\usepackage[utf8]{inputenc}
\usepackage[T1]{fontenc}
\usepackage[brazilian]{babel}
\usepackage{amsmath, amssymb, amsthm}
\usepackage{graphicx}
\usepackage{xcolor}
\usepackage{tikz}
\usetikzlibrary{arrows.meta, shapes, positioning, calc, decorations.pathreplacing, decorations.shapes, decorations.pathmorphing}
\usepackage{listings}
\usepackage{booktabs}
\usepackage{multirow}
\usepackage{hyperref}
\usepackage{chemfig}
\usepackage{chemarr}

% ============================================================================
% CONFIGURAÇÃO DO TEMA
% ============================================================================
\usetheme{Madrid}
\usecolortheme{default}

% Cores personalizadas
\definecolor{azulescuro}{RGB}{0, 51, 102}
\definecolor{azulclaro}{RGB}{51, 153, 255}
\definecolor{verdeacido}{RGB}{102, 204, 0}
\definecolor{laranjacel}{RGB}{255, 153, 51}
\definecolor{roxodna}{RGB}{153, 51, 255}

\setbeamercolor{structure}{fg=azulescuro}
\setbeamercolor{palette primary}{bg=azulescuro,fg=white}
\setbeamercolor{palette secondary}{bg=azulclaro,fg=white}
\setbeamercolor{palette tertiary}{bg=azulescuro,fg=white}
\setbeamercolor{palette quaternary}{bg=azulclaro,fg=white}
\setbeamercolor{block title}{bg=azulescuro,fg=white}
\setbeamercolor{block body}{bg=azulclaro!10,fg=black}
\setbeamercolor{block title alerted}{bg=red!80,fg=white}
\setbeamercolor{block body alerted}{bg=red!10,fg=black}
\setbeamercolor{block title example}{bg=verdeacido!80,fg=white}
\setbeamercolor{block body example}{bg=verdeacido!10,fg=black}

% ============================================================================
% CONFIGURAÇÃO DE CÓDIGO
% ============================================================================
\lstset{
    language=Python,
    basicstyle=\ttfamily\small,
    keywordstyle=\color{azulescuro}\bfseries,
    commentstyle=\color{verdeacido},
    stringstyle=\color{laranjacel},
    numbers=left,
    numberstyle=\tiny\color{gray},
    stepnumber=1,
    numbersep=5pt,
    backgroundcolor=\color{white},
    showspaces=false,
    showstringspaces=false,
    showtabs=false,
    frame=single,
    rulecolor=\color{black},
    tabsize=2,
    captionpos=b,
    breaklines=true,
    breakatwhitespace=false,
    escapeinside={\%*}{*)},
    xleftmargin=10pt,
    xrightmargin=10pt
}

% Definir estilo para R
\lstdefinestyle{Rstyle}{
    language=R,
    basicstyle=\ttfamily\small,
    keywordstyle=\color{azulescuro}\bfseries,
    commentstyle=\color{verdeacido},
    stringstyle=\color{laranjacel}
}

% Definir estilo para MATLAB
\lstdefinestyle{MATLABstyle}{
    language=Matlab,
    basicstyle=\ttfamily\small,
    keywordstyle=\color{azulescuro}\bfseries,
    commentstyle=\color{verdeacido},
    stringstyle=\color{laranjacel}
}

% ============================================================================
% COMANDOS PERSONALIZADOS
% ============================================================================

% Comando para equações destacadas
\newcommand{\destaque}[1]{%
    \begin{alertblock}{}
        \begin{center}
            #1
        \end{center}
    \end{alertblock}
}

% Comando para definições
\newcommand{\definicao}[2]{%
    \begin{block}{Definição: #1}
        #2
    \end{block}
}

% Comando para exemplos
\newcommand{\exemplo}[2]{%
    \begin{exampleblock}{Exemplo: #1}
        #2
    \end{exampleblock}
}

% Comando para observações importantes
\newcommand{\importante}[1]{%
    \begin{alertblock}{Importante!}
        #1
    \end{alertblock}
}

% Comandos matemáticos úteis
\newcommand{\R}{\mathbb{R}}
\newcommand{\N}{\mathbb{N}}
\newcommand{\Z}{\mathbb{Z}}
\newcommand{\dif}{\mathrm{d}}
\newcommand{\parcial}[2]{\frac{\partial #1}{\partial #2}}
\newcommand{\derivada}[2]{\frac{\dif #1}{\dif #2}}
\newcommand{\vect}[1]{\boldsymbol{#1}}

% Comandos biológicos
\newcommand{\gene}[1]{\textit{#1}}
\newcommand{\proteina}[1]{\textsc{#1}}
\newcommand{\especie}[1]{\textit{#1}}

% ============================================================================
% CONFIGURAÇÃO DE RODAPÉ
% ============================================================================
\setbeamertemplate{footline}{
  \leavevmode%
  \hbox{%
  \begin{beamercolorbox}[wd=.33\paperwidth,ht=2.25ex,dp=1ex,center]{author in head/foot}%
    \usebeamerfont{author in head/foot}\insertshortauthor
  \end{beamercolorbox}%
  \begin{beamercolorbox}[wd=.34\paperwidth,ht=2.25ex,dp=1ex,center]{title in head/foot}%
    \usebeamerfont{title in head/foot}\insertshorttitle
  \end{beamercolorbox}%
  \begin{beamercolorbox}[wd=.33\paperwidth,ht=2.25ex,dp=1ex,right]{date in head/foot}%
    \usebeamerfont{date in head/foot}\insertshortdate{}\hspace*{2em}
    \insertframenumber{} / \inserttotalframenumber\hspace*{2ex}
  \end{beamercolorbox}}%
  \vskip0pt%
}

% ============================================================================
% CONFIGURAÇÃO DE NAVEGAÇÃO
% ============================================================================
\setbeamertemplate{navigation symbols}{}

% ============================================================================
% CONFIGURAÇÃO DE ITEMIZE/ENUMERATE
% ============================================================================
\setbeamertemplate{itemize items}[circle]
\setbeamertemplate{enumerate items}[default]

% ============================================================================
% FIM DO TEMPLATE
% ============================================================================


% Pacotes adicionais para sequências
\usepackage{seqsplit}

% ============================================================================
% INFORMAÇÕES DO DOCUMENTO
% ============================================================================
\title[Cap. 6: Sistemas Gênicos]{Capítulo 6: Sistemas Gênicos}
\subtitle{Bioinformática para Biologia de Sistemas}
\author{Seu Nome}
\institute{Sua Instituição}
\date{\today}

\begin{document}

% ============================================================================
% SLIDE DE TÍTULO
% ============================================================================
\begin{frame}
\titlepage
\end{frame}

% ============================================================================
% SUMÁRIO
% ============================================================================
\begin{frame}{Sumário}
\tableofcontents
\end{frame}

% ============================================================================
% SEÇÃO 1: DOGMA CENTRAL
% ============================================================================
\section{Dogma Central da Biologia Molecular}

\begin{frame}{6.1 Dogma Central da Biologia Molecular}
\begin{center}
\begin{tikzpicture}[node distance=3.5cm, auto,
    box/.style={rectangle, draw, fill=azulclaro!20, text width=2.5cm, text centered, rounded corners, minimum height=1.2cm, font=\large}]

    \node[box, fill=roxodna!30] (dna) {\textbf{DNA}};
    \node[box, fill=laranjacel!30, right of=dna] (rna) {\textbf{RNA}};
    \node[box, fill=verdeacido!30, right of=rna] (prot) {\textbf{Proteína}};

    \draw[->, ultra thick, azulescuro, line width=2pt] (dna) -- node[above] {\textbf{Transcrição}} (rna);
    \draw[->, ultra thick, azulescuro, line width=2pt] (rna) -- node[above] {\textbf{Tradução}} (prot);
    \draw[->, ultra thick, azulescuro, line width=2pt] (dna) to [out=135,in=45,looseness=8] node[above] {\textbf{Replicação}} (dna);

\end{tikzpicture}
\end{center}

\vspace{0.5cm}

\importante{
O fluxo de informação genética segue a direção: DNA $\rightarrow$ RNA $\rightarrow$ Proteína
}

\begin{exampleblock}{Exceção: Retrovírus}
Transcriptase reversa: RNA $\rightarrow$ DNA
\end{exampleblock}
\end{frame}

\begin{frame}{Visão Geral dos Sistemas Gênicos}
\begin{block}{Objetivos de Aprendizagem}
\begin{itemize}
    \item Compreender a estrutura e organização dos genes
    \item Entender os mecanismos de regulação gênica
    \item Conhecer os diferentes tipos de RNA e suas funções
    \item Aprender técnicas de medição de expressão gênica
    \item Explorar bancos de dados genômicos (GenBank)
\end{itemize}
\end{block}

\vspace{0.3cm}

\definicao{Gene}{
Unidade fundamental da hereditariedade: sequência de DNA que codifica uma molécula funcional (RNA ou proteína) e inclui regiões regulatórias.
}
\end{frame}

% ============================================================================
% SEÇÃO 2: DNA FEATURES
% ============================================================================
\section{Características do DNA}

\begin{frame}{6.2 Características do DNA (DNA Features)}
\begin{columns}
\column{0.5\textwidth}
\textbf{Elementos Estruturais:}
\begin{itemize}
    \item Promotor
    \item Região codificante (ORF)
    \item Íntrons e éxons
    \item UTRs (5' e 3')
    \item Terminador
    \item Enhancers/Silencers
\end{itemize}

\column{0.5\textwidth}
\begin{center}
\begin{tikzpicture}[scale=0.7]
    % Gene structure
    \draw[thick] (0,0) -- (8,0);

    % Promoter
    \fill[roxodna!50] (0,0) rectangle (1,0.5);
    \node[above, font=\tiny] at (0.5,0.5) {Promotor};

    % Exons
    \fill[verdeacido!50] (1.5,0) rectangle (3,0.5);
    \fill[verdeacido!50] (4,0) rectangle (5.5,0.5);
    \node[above, font=\tiny] at (2.25,0.5) {Éxon 1};
    \node[above, font=\tiny] at (4.75,0.5) {Éxon 2};

    % Introns
    \draw[thick, azulclaro] (3,0.25) -- (4,0.25);
    \node[below, font=\tiny] at (3.5,-0.2) {Íntron};

    % Terminator
    \fill[roxodna!50] (7,0) rectangle (8,0.5);
    \node[above, font=\tiny] at (7.5,0.5) {Term};

    % Arrow
    \draw[->, thick] (0,-0.8) -- (8,-0.8);
    \node[below, font=\tiny] at (4,-0.8) {Direção da transcrição (5' $\rightarrow$ 3')};
\end{tikzpicture}
\end{center}
\end{columns}

\vspace{0.3cm}

\begin{exampleblock}{Procariotos vs Eucariotos}
\begin{itemize}
    \item \textbf{Procariotos:} genes sem íntrons, operons
    \item \textbf{Eucariotos:} íntrons, splicing alternativo
\end{itemize}
\end{exampleblock}
\end{frame}

\begin{frame}{Promotores e Elementos Regulatórios}
\begin{block}{Promotor}
Região onde a RNA polimerase se liga para iniciar a transcrição
\end{block}

\begin{columns}
\column{0.5\textwidth}
\textbf{Elementos em Procariotos:}
\begin{itemize}
    \item \textbf{-10 box} (Pribnow box):\\
    TATAAT
    \item \textbf{-35 box}:\\
    TTGACA
    \item \textbf{Operador:} controle negativo
\end{itemize}

\column{0.5\textwidth}
\textbf{Elementos em Eucariotos:}
\begin{itemize}
    \item \textbf{TATA box} (-25):\\
    TATAAA
    \item \textbf{CAAT box} (-75)
    \item \textbf{GC box}
    \item \textbf{Enhancers} (distantes)
\end{itemize}
\end{columns}

\vspace{0.3cm}

\destaque{
\texttt{5'---[Enhancer]---[CAAT]---[TATA]---[+1]---Gene---3'}
}
\end{frame}

\begin{frame}{Open Reading Frame (ORF)}
\definicao{ORF}{
Sequência contínua de códons entre um códon de início (ATG) e um códon de parada (TAA, TAG, TGA), sem códons de parada internos.
}

\vspace{0.3cm}

\begin{exampleblock}{Exemplo}
\texttt{5'-ATG GCT AAA TGG CCA TAA-3'}\\
\texttt{~~~~Met-Ala-Lys-Trp-Pro-STOP}

\vspace{0.2cm}

\textbf{ORF:} 15 nucleotídeos = 5 códons
\end{exampleblock}

\begin{block}{Identificação de ORFs}
Critérios:
\begin{itemize}
    \item Começa com ATG (códon de início)
    \item Termina com TAA, TAG ou TGA (códons de parada)
    \item Tamanho mínimo (geralmente \textgreater 300 nt / 100 aa)
    \item Mesma fase de leitura (frame)
\end{itemize}
\end{block}
\end{frame}

\begin{frame}[fragile]{Busca de ORFs Computacional}
\begin{lstlisting}[language=Python, caption=Identificação de ORFs]
def find_orfs(sequence, min_length=100):
    """
    Encontra ORFs em uma sequencia de DNA
    """
    start_codon = "ATG"
    stop_codons = ["TAA", "TAG", "TGA"]
    orfs = []

    # Para cada frame (0, 1, 2)
    for frame in range(3):
        for i in range(frame, len(sequence)-2, 3):
            codon = sequence[i:i+3]

            # Encontrou inicio
            if codon == start_codon:
                # Procurar stop codon
                for j in range(i+3, len(sequence)-2, 3):
                    stop = sequence[j:j+3]
                    if stop in stop_codons:
                        orf_len = j - i + 3
                        if orf_len >= min_length*3:
                            orfs.append({
                                'start': i,
                                'end': j+3,
                                'length': orf_len,
                                'frame': frame
                            })
                        break
    return orfs
\end{lstlisting}
\end{frame}

% ============================================================================
% SEÇÃO 3: REGULAÇÃO GÊNICA
% ============================================================================
\section{Regulação Gênica}

\begin{frame}{6.3 Regulação Gênica}
\begin{block}{Níveis de Regulação}
\begin{enumerate}
    \item \textbf{Transcricional:} controle da síntese de RNA
    \item \textbf{Pós-transcricional:} splicing, estabilidade do mRNA
    \item \textbf{Traducional:} controle da síntese proteica
    \item \textbf{Pós-traducional:} modificações, degradação
    \item \textbf{Epigenética:} metilação, acetilação de histonas
\end{enumerate}
\end{block}

\vspace{0.3cm}

\importante{
A regulação da expressão gênica permite:
\begin{itemize}
    \item Resposta a sinais ambientais
    \item Diferenciação celular
    \item Desenvolvimento embrionário
    \item Homeostase
\end{itemize}
}
\end{frame}

\begin{frame}{Regulação Transcricional em Procariotos}
\exemplo{Operon lac (Lactose)}{
Sistema clássico de regulação negativa e positiva
}

\begin{center}
\begin{tikzpicture}[scale=0.7, every node/.style={font=\tiny}]
    % Gene structure
    \draw[thick] (0,0) -- (10,0);

    % Promoter
    \fill[roxodna!50] (0.5,0) rectangle (1.5,0.5);
    \node[above] at (1,0.5) {P};

    % Operator
    \fill[red!50] (1.5,0) rectangle (2.5,0.5);
    \node[above] at (2,0.5) {O};

    % Genes
    \fill[verdeacido!50] (3,0) rectangle (4.5,0.5);
    \node[above] at (3.75,0.5) {\gene{lacZ}};

    \fill[verdeacido!50] (5,0) rectangle (6.5,0.5);
    \node[above] at (5.75,0.5) {\gene{lacY}};

    \fill[verdeacido!50] (7,0) rectangle (8.5,0.5);
    \node[above] at (7.75,0.5) {\gene{lacA}};

    % Repressor (absent)
    \node[draw, circle, fill=yellow!50, minimum size=0.8cm] at (2,-1.5) {LacI};
    \draw[->, thick, red] (2,-1) -- (2,-0.5);

    % Lactose
    \node[draw, circle, fill=azulclaro!50] at (0,-1.5) {Lac};
    \draw[->, dashed] (0.5,-1.5) -- (1.5,-1.5);

    \node[below] at (5,-2) {Sem lactose: repressor ativo (OFF)};
\end{tikzpicture}
\end{center}

\begin{itemize}
    \item \textbf{Sem lactose:} repressor LacI bloqueia transcrição
    \item \textbf{Com lactose:} lactose liga a LacI, liberando operador
\end{itemize}
\end{frame}

\begin{frame}{Regulação Transcricional em Eucariotos}
\begin{columns}
\column{0.5\textwidth}
\textbf{Fatores de Transcrição:}
\begin{itemize}
    \item Ativadores
    \item Repressores
    \item Mediadores
    \item Remodeladores de cromatina
\end{itemize}

\vspace{0.3cm}

\textbf{Estrutura da Cromatina:}
\begin{itemize}
    \item Eucromatina (aberta)
    \item Heterocromatina (fechada)
    \item Nucleossomos
\end{itemize}

\column{0.5\textwidth}
\begin{center}
\begin{tikzpicture}[scale=0.8]
    % DNA wrapping around histone
    \draw[thick, roxodna, domain=0:360, samples=100, smooth]
        plot ({1.5 + 0.8*cos(\x)}, {0.02*\x + 0.8*sin(\x)});

    % Histone core
    \fill[azulclaro!50] (1.5,1.5) circle (0.7);
    \node at (1.5,1.5) {H};

    % Labels
    \node[below, font=\tiny] at (1.5,0) {Nucleossomo};

    % Modificações
    \foreach \angle in {45, 135, 225, 315} {
        \fill[laranjacel] ({1.5 + 0.9*cos(\angle)}, {1.5 + 0.9*sin(\angle)}) circle (1.5pt);
    }
    \node[right, font=\tiny] at (2.5,1.5) {Modificações};
\end{tikzpicture}
\end{center}
\end{columns}

\vspace{0.3cm}

\begin{exampleblock}{Modificações Epigenéticas}
\begin{itemize}
    \item Acetilação: ativa transcrição
    \item Metilação: silencia genes
\end{itemize}
\end{exampleblock}
\end{frame}

\begin{frame}{Modelo Matemático de Regulação Gênica}
\begin{block}{Equação de Hill}
Modelo para regulação cooperativa:
\begin{equation*}
\frac{d[mRNA]}{dt} = \frac{V_{max} [TF]^n}{K_d^n + [TF]^n} - \delta [mRNA]
\end{equation*}

onde:
\begin{itemize}
    \item $[TF]$ = concentração do fator de transcrição
    \item $n$ = coeficiente de Hill (cooperatividade)
    \item $K_d$ = constante de dissociação
    \item $\delta$ = taxa de degradação do mRNA
\end{itemize}
\end{block}

\begin{exampleblock}{Interpretação de $n$}
\begin{itemize}
    \item $n = 1$: sem cooperatividade (Michaelis-Menten)
    \item $n > 1$: cooperatividade positiva (switch-like)
    \item $n < 1$: cooperatividade negativa
\end{itemize}
\end{exampleblock}
\end{frame}

\begin{frame}[fragile]{Simulação de Regulação Gênica}
\begin{lstlisting}[language=Python, caption=Função de Hill]
import numpy as np
import matplotlib.pyplot as plt

def hill_function(TF, Vmax, Kd, n):
    """Funcao de Hill para regulacao genica"""
    return Vmax * TF**n / (Kd**n + TF**n)

# Parametros
Vmax = 100
Kd = 10
TF = np.linspace(0, 50, 200)

# Plotar diferentes cooperatividades
plt.figure(figsize=(10, 6))
for n in [1, 2, 4, 8]:
    mRNA = hill_function(TF, Vmax, Kd, n)
    plt.plot(TF, mRNA, linewidth=2, label=f'n = {n}')

plt.axhline(y=Vmax/2, color='gray', linestyle='--')
plt.axvline(x=Kd, color='gray', linestyle='--')
plt.xlabel('[Fator de Transcricao]')
plt.ylabel('Taxa de Transcricao')
plt.title('Regulacao Genica: Efeito da Cooperatividade')
plt.legend()
plt.grid(True)
\end{lstlisting}
\end{frame}

% ============================================================================
% SEÇÃO 4: TIPOS DE RNA
% ============================================================================
\section{Tipos de RNA}

\begin{frame}{6.4-6.6 Tipos de RNA}
\begin{block}{Classificação Funcional}
\begin{columns}
\column{0.5\textwidth}
\textbf{RNAs Codificantes:}
\begin{itemize}
    \item \textbf{mRNA:} mensageiro\\
    (proteínas)
\end{itemize}

\vspace{0.3cm}

\textbf{RNAs Não-codificantes:}
\begin{itemize}
    \item \textbf{rRNA:} ribossomal\\
    (estrutura, catálise)
    \item \textbf{tRNA:} transferência\\
    (adaptador)
\end{itemize}

\column{0.5\textwidth}
\textbf{RNAs Regulatórios:}
\begin{itemize}
    \item \textbf{miRNA:} microRNA\\
    (silenciamento)
    \item \textbf{siRNA:} RNA interferente\\
    (degradação)
    \item \textbf{lncRNA:} longo não-codificante\\
    (regulação cromossômica)
    \item \textbf{snoRNA:} nucleolar pequeno\\
    (modificação de rRNA)
\end{itemize}
\end{columns}
\end{block}
\end{frame}

\begin{frame}{6.4 RNA Mensageiro (mRNA)}
\begin{block}{Estrutura do mRNA em Eucariotos}
\begin{center}
\begin{tikzpicture}[scale=0.8]
    % mRNA structure
    \draw[thick, laranjacel] (0,0) -- (10,0);

    % 5' cap
    \node[draw, circle, fill=roxodna!50, minimum size=0.6cm] at (0,0) {Cap};
    \node[above, font=\tiny] at (0,0.5) {5' m$^7$G};

    % 5' UTR
    \draw[thick] (0.5,0) -- (2,0);
    \node[above, font=\tiny] at (1.25,0.3) {5' UTR};

    % Start codon
    \fill[verdeacido!50] (2,0) circle (2pt);
    \node[below, font=\tiny] at (2,-0.3) {AUG};

    % Coding region
    \fill[verdeacido!50] (2,-0.2) rectangle (7,0.2);
    \node[above, font=\tiny] at (4.5,0.3) {Região Codificante};

    % Stop codon
    \fill[red!50] (7,0) circle (2pt);
    \node[below, font=\tiny] at (7,-0.3) {STOP};

    % 3' UTR
    \draw[thick] (7,0) -- (9,0);
    \node[above, font=\tiny] at (8,0.3) {3' UTR};

    % Poly-A tail
    \draw[thick, azulclaro] (9,0) -- (10,0);
    \node[above, font=\tiny] at (9.5,0.3) {Poly(A)};
    \node[below, font=\tiny] at (9.5,-0.3) {AAAA...};
\end{tikzpicture}
\end{center}
\end{block}

\textbf{Elementos Funcionais:}
\begin{itemize}
    \item \textbf{5' Cap:} proteção, reconhecimento para tradução
    \item \textbf{5' UTR:} regulação traducional
    \item \textbf{CDS:} sequência codificante (Coding Sequence)
    \item \textbf{3' UTR:} estabilidade, localização
    \item \textbf{Poly(A):} estabilidade, exportação nuclear
\end{itemize}
\end{frame}

\begin{frame}{Processamento do mRNA em Eucariotos}
\begin{center}
\begin{tikzpicture}[node distance=2cm, auto,
    box/.style={rectangle, draw, text width=3cm, text centered, rounded corners, minimum height=0.8cm, font=\small}]

    \node[box, fill=roxodna!20] (pre) {pré-mRNA\\(transcrito primário)};
    \node[box, fill=laranjacel!20, below of=pre] (splice) {Splicing};
    \node[box, fill=azulclaro!20, below of=splice] (mature) {mRNA maduro};
    \node[box, fill=verdeacido!20, below of=mature] (export) {Exportação};

    \draw[->, thick] (pre) -- (splice);
    \draw[->, thick] (splice) -- (mature);
    \draw[->, thick] (mature) -- (export);

    % Annotations
    \node[right, font=\tiny, text width=3cm] at (4,0) {Éxons + Íntrons\\+ 5' Cap + Poly(A)};
    \node[right, font=\tiny, text width=3cm] at (4,-2) {Remoção de íntrons\\Junção de éxons};
    \node[right, font=\tiny, text width=3cm] at (4,-4) {Apenas éxons\\+ Cap + Poly(A)};
    \node[right, font=\tiny, text width=3cm] at (4,-6) {Núcleo $\rightarrow$ Citoplasma};
\end{tikzpicture}
\end{center}
\end{frame}

\begin{frame}{6.5 RNA de Transferência (tRNA)}
\begin{columns}
\column{0.5\textwidth}
\textbf{Estrutura em Folha de Trevo:}
\begin{itemize}
    \item Braço aceptor (3' CCA)
    \item Braço D
    \item Braço anticódon
    \item Braço T$\Psi$C
    \item Braço variável
\end{itemize}

\vspace{0.3cm}

\textbf{Função:}
\begin{itemize}
    \item Transporta aminoácidos
    \item Reconhece códons via anticódon
    \item Adaptador molecular
\end{itemize}

\column{0.5\textwidth}
\begin{center}
\begin{tikzpicture}[scale=0.6]
    % tRNA cloverleaf
    % Acceptor stem
    \draw[thick, azulescuro] (0,0) -- (0,2);
    \draw[thick, azulescuro] (0.8,0) -- (0.8,2);
    \foreach \y in {0.3, 0.7, 1.1, 1.5}
        \draw (0,\y) -- (0.8,\y);

    % 3' end
    \node[above, font=\tiny] at (0.4,2) {3' CCA-AA};

    % D arm
    \draw[thick, azulescuro] (-2,1) circle (0.6);
    \node[font=\tiny] at (-2,1) {D};

    % Anticodon arm
    \draw[thick, azulescuro] (0,-1) circle (0.8);
    \node[font=\tiny] at (0,-0.5) {Anticódon};
    \node[below, font=\scriptsize, fill=laranjacel!50] at (0,-1.8) {UAC};

    % TψC arm
    \draw[thick, azulescuro] (2.8,1) circle (0.6);
    \node[font=\tiny] at (2.8,1) {T$\Psi$C};

    % Connections
    \draw[thick, azulescuro] (0,0.5) -- (-1.5,1);
    \draw[thick, azulescuro] (0.8,0.5) -- (2.3,1);
    \draw[thick, azulescuro] (0,-0.2) -- (0,-0.2);
\end{tikzpicture}
\end{center}

\begin{exampleblock}{Wobble Pairing}
Flexibilidade na 3ª posição do códon
\end{exampleblock}
\end{columns}
\end{frame}

\begin{frame}{6.6 RNA Ribossomal (rRNA)}
\begin{block}{Componente Estrutural dos Ribossomos}
\begin{center}
\begin{tabular}{lccc}
\toprule
\textbf{Organismo} & \textbf{Subunidade} & \textbf{rRNA} & \textbf{Proteínas} \\
\midrule
\multirow{2}{*}{Procarioto} & 30S & 16S & 21 \\
                             & 50S & 23S, 5S & 31 \\
\midrule
\multirow{2}{*}{Eucarioto} & 40S & 18S & 33 \\
                           & 60S & 28S, 5.8S, 5S & 49 \\
\bottomrule
\end{tabular}
\end{center}
\end{block}

\vspace{0.3cm}

\begin{columns}
\column{0.5\textwidth}
\textbf{Funções:}
\begin{itemize}
    \item Estrutura do ribossomo
    \item Atividade peptidil-transferase
    \item Ligação de tRNA
    \item Translocação
\end{itemize}

\column{0.5\textwidth}
\importante{
rRNA é uma \textbf{ribozima}: RNA com atividade catalítica!
}
\end{columns}
\end{frame}

\begin{frame}{RNAs Regulatórios: microRNA (miRNA)}
\begin{block}{Biogênese do miRNA}
\begin{enumerate}
    \item \textbf{Transcrição:} pri-miRNA (longo, hairpin)
    \item \textbf{Processamento nuclear:} Drosha $\rightarrow$ pre-miRNA ($\sim$70 nt)
    \item \textbf{Exportação:} Exportina-5
    \item \textbf{Processamento citoplasmático:} Dicer $\rightarrow$ miRNA maduro ($\sim$22 nt)
    \item \textbf{RISC:} complexo de silenciamento
\end{enumerate}
\end{block}

\vspace{0.3cm}

\begin{exampleblock}{Mecanismo de Ação}
\begin{itemize}
    \item \textbf{Pareamento perfeito:} clivagem do mRNA-alvo
    \item \textbf{Pareamento imperfeito:} repressão traducional
    \item Um miRNA pode regular centenas de mRNAs
\end{itemize}
\end{exampleblock}
\end{frame}

% ============================================================================
% SEÇÃO 5: MEDIÇÃO DE EXPRESSÃO GÊNICA
% ============================================================================
\section{Medição de Expressão Gênica}

\begin{frame}{6.7 Medição de Expressão Gênica}
\begin{block}{Técnicas Principais}
\begin{columns}
\column{0.5\textwidth}
\textbf{Nível de mRNA:}
\begin{itemize}
    \item Northern blot
    \item RT-PCR / qPCR
    \item Microarrays
    \item RNA-seq
    \item Single-cell RNA-seq
\end{itemize}

\column{0.5\textwidth}
\textbf{Nível de Proteína:}
\begin{itemize}
    \item Western blot
    \item Imunofluorescência
    \item Espectrometria de massa
    \item Flow cytometry
\end{itemize}
\end{columns}
\end{block}

\vspace{0.3cm}

\importante{
mRNA $\not\propto$ Proteína sempre!\\
Regulação pós-transcricional e pós-traducional
}
\end{frame}

\begin{frame}{Microarrays de DNA}
\begin{columns}
\column{0.5\textwidth}
\textbf{Princípio:}
\begin{enumerate}
    \item Sondas fixadas em chip
    \item Hibridização de cDNA marcado
    \item Detecção fluorescente
    \item Análise de intensidade
\end{enumerate}

\vspace{0.3cm}

\textbf{Aplicações:}
\begin{itemize}
    \item Perfil de expressão global
    \item Comparação de condições
    \item Classificação de amostras
    \item Descoberta de genes
\end{itemize}

\column{0.5\textwidth}
\begin{center}
\begin{tikzpicture}[scale=0.7]
    % Microarray chip
    \draw[thick] (0,0) rectangle (4,4);

    % Grid of spots
    \foreach \x in {0.5, 1.0, 1.5, 2.0, 2.5, 3.0, 3.5} {
        \foreach \y in {0.5, 1.0, 1.5, 2.0, 2.5, 3.0, 3.5} {
            \pgfmathsetmacro{\intensity}{random(0,100)}
            \definecolor{spotcolor}{rgb}{1,\intensity/100,0}
            \fill[spotcolor] (\x,\y) circle (1.5pt);
        }
    }

    \node[below] at (2,-0.3) {Microarray};

    % Legend
    \draw[fill=red] (4.5,3.5) rectangle (5,4);
    \node[right, font=\tiny] at (5.1,3.75) {Alta expressão};

    \draw[fill=yellow] (4.5,2.5) rectangle (5,3);
    \node[right, font=\tiny] at (5.1,2.75) {Média};

    \draw[fill=orange] (4.5,1.5) rectangle (5,2);
    \node[right, font=\tiny] at (5.1,1.75) {Baixa};
\end{tikzpicture}
\end{center}
\end{columns}
\end{frame}

\begin{frame}{RNA-seq: Sequenciamento de RNA}
\begin{block}{Vantagens sobre Microarrays}
\begin{itemize}
    \item Maior faixa dinâmica
    \item Sem necessidade de sondas pré-definidas
    \item Detecção de novos transcritos
    \item Resolução de nucleotídeo único
    \item Identificação de variantes de splicing
    \item Quantificação absoluta (não relativa)
\end{itemize}
\end{block}

\begin{exampleblock}{Pipeline RNA-seq}
\begin{enumerate}
    \item Extração de RNA
    \item Preparação de biblioteca (cDNA)
    \item Sequenciamento (Illumina, PacBio, etc.)
    \item Alinhamento a genoma de referência
    \item Contagem de reads por gene
    \item Análise de expressão diferencial
\end{enumerate}
\end{exampleblock}
\end{frame}

\begin{frame}[fragile]{Análise de Dados de RNA-seq}
\begin{lstlisting}[language=Python, caption=Análise de expressão diferencial]
import pandas as pd
import numpy as np
from scipy import stats

# Carregar contagens (genes x amostras)
counts = pd.read_csv('gene_counts.csv', index_col=0)

# Normalizar (TPM, RPKM, ou DESeq2)
def tpm_normalize(counts, gene_lengths):
    """Normalize para TPM (Transcripts Per Million)"""
    # Reads per kilobase
    rpk = counts / (gene_lengths / 1000)
    # Scaling factor
    scaling_factor = rpk.sum() / 1e6
    tpm = rpk / scaling_factor
    return tpm

# Teste t para expressao diferencial
def differential_expression(control, treatment):
    """Teste t para cada gene"""
    t_stat, p_value = stats.ttest_ind(
        control, treatment, axis=1
    )
    log_fc = np.log2(
        treatment.mean(axis=1) / control.mean(axis=1)
    )
    return log_fc, p_value

# Aplicar
log_fc, p_val = differential_expression(
    counts.iloc[:, :3],  # controle
    counts.iloc[:, 3:]   # tratamento
)
\end{lstlisting}
\end{frame}

\begin{frame}[fragile]{Visualização: Volcano Plot}
\begin{lstlisting}[language=Python, caption=Volcano plot para DE genes]
import matplotlib.pyplot as plt

# Calcular -log10(p-value)
neg_log_p = -np.log10(p_val)

# Criar volcano plot
plt.figure(figsize=(10, 8))
plt.scatter(log_fc, neg_log_p, alpha=0.5, s=10)

# Threshold lines
plt.axhline(y=-np.log10(0.05), color='r',
            linestyle='--', label='p = 0.05')
plt.axvline(x=-1, color='b', linestyle='--')
plt.axvline(x=1, color='b', linestyle='--',
            label='|log2 FC| = 1')

# Destacar genes significativos
sig = (p_val < 0.05) & (np.abs(log_fc) > 1)
plt.scatter(log_fc[sig], neg_log_p[sig],
            color='red', s=20, label='Significativo')

plt.xlabel('log2(Fold Change)')
plt.ylabel('-log10(p-value)')
plt.title('Volcano Plot: Expressao Diferencial')
plt.legend()
plt.grid(True, alpha=0.3)
\end{lstlisting}
\end{frame}

% ============================================================================
% SEÇÃO 6: LOCALIZAÇÃO DE GENES
% ============================================================================
\section{Localização de Genes}

\begin{frame}{6.8 Localização de Genes (Gene Localization)}
\begin{block}{Predição de Genes \textit{Ab Initio}}
Métodos computacionais que identificam genes baseando-se apenas na sequência de DNA:

\begin{itemize}
    \item Busca de ORFs
    \item Detecção de sinais (promotores, splice sites)
    \item Modelos de Markov Ocultos (HMM)
    \item Redes neurais
\end{itemize}
\end{block}

\begin{block}{Predição Baseada em Evidências}
Usa dados experimentais:
\begin{itemize}
    \item Alinhamento de ESTs (Expressed Sequence Tags)
    \item Alinhamento de proteínas homólogas
    \item Dados de RNA-seq
    \item Sintenia (comparação genômica)
\end{itemize}
\end{block}
\end{frame}

\begin{frame}{Ferramentas de Predição de Genes}
\begin{columns}
\column{0.5\textwidth}
\textbf{Procariotos:}
\begin{itemize}
    \item \textbf{Prodigal}
    \item \textbf{GeneMark}
    \item \textbf{Glimmer}
    \item \textbf{RAST}
\end{itemize}

\vspace{0.3cm}

\begin{exampleblock}{Características}
\begin{itemize}
    \item Genes densos
    \item Poucos íntrons
    \item Operons
    \item Mais simples
\end{itemize}
\end{exampleblock}

\column{0.5\textwidth}
\textbf{Eucariotos:}
\begin{itemize}
    \item \textbf{Augustus}
    \item \textbf{GenScan}
    \item \textbf{SNAP}
    \item \textbf{MAKER} (pipeline)
    \item \textbf{Braker}
\end{itemize}

\vspace{0.3cm}

\begin{alertblock}{Desafios}
\begin{itemize}
    \item Íntrons
    \item Splicing alternativo
    \item Genes esparsos
    \item DNA repetitivo
\end{itemize}
\end{alertblock}
\end{columns}
\end{frame}

\begin{frame}[fragile]{Exemplo: Usando BioPython para Anotação}
\begin{lstlisting}[language=Python, caption=Busca de ORFs com BioPython]
from Bio import SeqIO
from Bio.SeqUtils import GC

def annotate_sequence(fasta_file):
    """Anotacao basica de sequencia"""
    for record in SeqIO.parse(fasta_file, "fasta"):
        seq = record.seq
        print(f"ID: {record.id}")
        print(f"Comprimento: {len(seq)} bp")
        print(f"GC%: {GC(seq):.2f}%")

        # Buscar ORFs em todas as 6 frames
        for strand, nuc in [(+1, seq), (-1, seq.reverse_complement())]:
            for frame in range(3):
                length = 3 * ((len(seq) - frame) // 3)
                for pro in nuc[frame:frame+length].translate().split("*"):
                    if len(pro) >= 100:  # min 100 aa
                        print(f"Frame {strand}{frame}: ORF de {len(pro)} aa")

# Usar
annotate_sequence("genome.fasta")
\end{lstlisting}
\end{frame}

% ============================================================================
% SEÇÃO 7: GENBANK
% ============================================================================
\section{GenBank}

\begin{frame}{GenBank: Banco de Dados de Sequências}
\begin{block}{O que é GenBank?}
\begin{itemize}
    \item Banco de dados público do NCBI (National Center for Biotechnology Information)
    \item Coleção abrangente de sequências de DNA
    \item Anotações de genes, proteínas, características
    \item Atualizado diariamente
    \item Livre acesso
\end{itemize}
\end{block}

\vspace{0.3cm}

\begin{exampleblock}{Estatísticas (2024)}
\begin{itemize}
    \item \textgreater 350 milhões de sequências
    \item \textgreater 10 trilhões de pares de bases
    \item Crescimento exponencial
\end{itemize}
\end{exampleblock}

\begin{block}{URL}
\url{https://www.ncbi.nlm.nih.gov/genbank/}
\end{block}
\end{frame}

\begin{frame}{Formato de Arquivo GenBank}
\begin{block}{Estrutura}
\begin{verbatim}
LOCUS       NM_000518     5510 bp    mRNA
DEFINITION  Homo sapiens hemoglobin beta
ACCESSION   NM_000518
VERSION     NM_000518.5
ORGANISM    Homo sapiens
FEATURES             Location/Qualifiers
     gene            1..5510
                     /gene="HBB"
     CDS             51..495
                     /product="hemoglobin beta"
                     /translation="MVHLTPEEKSAVTAL..."
ORIGIN
        1 acatttgctt ctgacacaac tgtgttcact agcaacctca
       41 aacagacacc atggtgcatc tgactcctga ggagaagtct
       ...
//
\end{verbatim}
\end{block}
\end{frame}

\begin{frame}[fragile]{Acessando GenBank com BioPython}
\begin{lstlisting}[language=Python, caption=Download de sequências]
from Bio import Entrez, SeqIO

# Sempre fornecer email
Entrez.email = "seu_email@example.com"

# Buscar por gene
def search_genbank(query, max_results=10):
    """Busca no GenBank"""
    handle = Entrez.esearch(
        db="nucleotide",
        term=query,
        retmax=max_results
    )
    record = Entrez.read(handle)
    return record["IdList"]

# Download de sequencia
def fetch_sequence(genbank_id):
    """Baixa sequencia do GenBank"""
    handle = Entrez.efetch(
        db="nucleotide",
        id=genbank_id,
        rettype="gb",
        retmode="text"
    )
    record = SeqIO.read(handle, "genbank")
    return record

# Exemplo
ids = search_genbank("Homo sapiens HBB")
seq = fetch_sequence(ids[0])
print(f"Gene: {seq.annotations['organism']}")
print(f"Tamanho: {len(seq)} bp")
\end{lstlisting}
\end{frame}

\begin{frame}[fragile]{Análise de Features no GenBank}
\begin{lstlisting}[language=Python, caption=Extrair features]
def extract_features(record):
    """Extrai features de um registro GenBank"""
    features_dict = {
        'genes': [],
        'CDS': [],
        'mRNA': []
    }

    for feature in record.features:
        if feature.type == 'gene':
            gene_name = feature.qualifiers.get('gene', ['Unknown'])[0]
            location = f"{feature.location.start}-{feature.location.end}"
            features_dict['genes'].append({
                'name': gene_name,
                'location': location
            })

        elif feature.type == 'CDS':
            product = feature.qualifiers.get('product', ['Unknown'])[0]
            translation = feature.qualifiers.get('translation', [''])[0]
            features_dict['CDS'].append({
                'product': product,
                'length': len(translation)
            })

    return features_dict

# Usar
features = extract_features(seq)
print(f"Genes encontrados: {len(features['genes'])}")
print(f"CDSs encontradas: {len(features['CDS'])}")
\end{lstlisting}
\end{frame}

% ============================================================================
% SEÇÃO 8: ANÁLISE INTEGRADA
% ============================================================================
\section{Análise Integrada}

\begin{frame}{Análise de Sequências: Pipeline Completo}
\begin{center}
\begin{tikzpicture}[node distance=1.5cm, auto,
    box/.style={rectangle, draw, text width=3cm, text centered, rounded corners, minimum height=0.8cm, font=\small}]

    \node[box, fill=roxodna!30] (seq) {Sequência\\Genômica};
    \node[box, fill=laranjacel!30, below of=seq] (pred) {Predição de\\Genes};
    \node[box, fill=verdeacido!30, below of=pred] (annot) {Anotação\\Funcional};
    \node[box, fill=azulclaro!30, below of=annot] (expr) {Análise de\\Expressão};
    \node[box, fill=yellow!30, below of=expr] (net) {Redes\\Regulatórias};

    \draw[->, thick] (seq) -- (pred);
    \draw[->, thick] (pred) -- (annot);
    \draw[->, thick] (annot) -- (expr);
    \draw[->, thick] (expr) -- (net);

    % Tools
    \node[right, font=\tiny, text width=3cm] at (4,0) {FASTA/GenBank};
    \node[right, font=\tiny, text width=3cm] at (4,-1.5) {Augustus, Prodigal};
    \node[right, font=\tiny, text width=3cm] at (4,-3) {BLAST, InterPro};
    \node[right, font=\tiny, text width=3cm] at (4,-4.5) {RNA-seq, qPCR};
    \node[right, font=\tiny, text width=3cm] at (4,-6) {Cytoscape, STRING};
\end{tikzpicture}
\end{center}
\end{frame}

\begin{frame}{Bancos de Dados Relacionados}
\begin{columns}
\column{0.5\textwidth}
\begin{block}{NCBI}
\begin{itemize}
    \item \textbf{GenBank:} sequências
    \item \textbf{RefSeq:} curadas
    \item \textbf{GEO:} expressão
    \item \textbf{SRA:} sequenciamento
    \item \textbf{PubMed:} literatura
\end{itemize}
\end{block}

\begin{block}{EBI (Europa)}
\begin{itemize}
    \item \textbf{ENA:} sequências
    \item \textbf{ArrayExpress:} microarrays
    \item \textbf{Ensembl:} genomas
\end{itemize}
\end{block}

\column{0.5\textwidth}
\begin{block}{Especializados}
\begin{itemize}
    \item \textbf{KEGG:} vias metabólicas
    \item \textbf{GO:} ontologia
    \item \textbf{UniProt:} proteínas
    \item \textbf{miRBase:} microRNAs
    \item \textbf{Rfam:} RNAs não-codificantes
\end{itemize}
\end{block}

\begin{exampleblock}{Integração}
Ferramentas como DAVID, Enrichr e WebGestalt integram múltiplos bancos
\end{exampleblock}
\end{columns}
\end{frame}

\begin{frame}{Ontologia Gênica (GO)}
\definicao{Gene Ontology}{
Vocabulário controlado para descrever funções de genes em três domínios:
\begin{itemize}
    \item \textbf{MF:} Molecular Function (função molecular)
    \item \textbf{BP:} Biological Process (processo biológico)
    \item \textbf{CC:} Cellular Component (componente celular)
\end{itemize}
}

\vspace{0.3cm}

\begin{exampleblock}{Exemplo: Hemoglobina}
\begin{itemize}
    \item \textbf{MF:} oxygen transporter activity (GO:0005344)
    \item \textbf{BP:} oxygen transport (GO:0015671)
    \item \textbf{CC:} hemoglobin complex (GO:0005833)
\end{itemize}
\end{exampleblock}

\begin{alertblock}{Análise de Enriquecimento}
Identificar processos/funções super-representados em conjunto de genes
\end{alertblock}
\end{frame}

% ============================================================================
% SEÇÃO 9: ESTUDOS DE CASO
% ============================================================================
\section{Estudos de Caso}

\begin{frame}{Caso 1: Análise de Expressão em Câncer}
\begin{block}{Objetivo}
Identificar genes diferencialmente expressos entre tecido normal e tumoral
\end{block}

\textbf{Abordagem:}
\begin{enumerate}
    \item Coletar dados de RNA-seq (GEO/TCGA)
    \item Normalizar e filtrar
    \item Análise de expressão diferencial (DESeq2)
    \item Correção para múltiplos testes (FDR)
    \item Análise de enriquecimento (GO, KEGG)
    \item Validação experimental (qPCR)
\end{enumerate}

\vspace{0.3cm}

\begin{exampleblock}{Resultado Esperado}
\begin{itemize}
    \item Lista de genes super/sub-expressos
    \item Vias biológicas alteradas
    \item Potenciais biomarcadores
    \item Alvos terapêuticos
\end{itemize}
\end{exampleblock}
\end{frame}

\begin{frame}{Caso 2: Descoberta de microRNAs Regulatórios}
\begin{block}{Objetivo}
Identificar miRNAs que regulam genes envolvidos em resposta ao estresse
\end{block}

\textbf{Pipeline:}
\begin{enumerate}
    \item Perfil de expressão de miRNAs (small RNA-seq)
    \item Perfil de expressão de mRNAs (RNA-seq)
    \item Predição de alvos (TargetScan, miRanda)
    \item Correlação inversa miRNA-mRNA
    \item Validação experimental (luciferase assay)
\end{enumerate}

\vspace{0.3cm}

\begin{alertblock}{Ferramentas}
\begin{itemize}
    \item \textbf{miRBase:} banco de miRNAs
    \item \textbf{TargetScan:} predição de alvos
    \item \textbf{DIANA-TarBase:} interações validadas
\end{itemize}
\end{alertblock}
\end{frame}

% ============================================================================
% SEÇÃO 10: EXERCÍCIOS
% ============================================================================
\section{Exercícios}

\begin{frame}{Exercícios - Parte 1}
\begin{block}{Questão 1: Cálculo de Códons}
Uma região codificante tem 1200 nucleotídeos.
\begin{enumerate}
    \item Quantos códons ela contém?
    \item Quantos aminoácidos na proteína resultante?
    \item Se a taxa de tradução é 15 aa/segundo, quanto tempo leva?
\end{enumerate}
\end{block}

\begin{block}{Questão 2: Conteúdo GC}
Por que o conteúdo GC é importante? Calcule o conteúdo GC da sequência:\\
\texttt{5'-ATGCGATCGTAGCTAGCTA-3'}
\end{block}

\begin{block}{Questão 3: Função de Hill}
Para uma função de Hill com $n=4$, $K_d=10$, $V_{max}=100$:
\begin{enumerate}
    \item Calcule a expressão quando $[TF]=10$
    \item Em que concentração se atinge 90\% de $V_{max}$?
\end{enumerate}
\end{block}
\end{frame}

\begin{frame}{Exercícios - Parte 2}
\begin{alertblock}{Questão 4: Projeto Prático - Análise de Sequência}
Escolha um gene de interesse (ex: TP53, BRCA1, MYC):
\begin{enumerate}
    \item Baixe a sequência do GenBank usando BioPython
    \item Identifique e anote todas as features
    \item Calcule estatísticas básicas (tamanho, GC\%, nº de éxons)
    \item Identifique potenciais sítios de ligação de fatores de transcrição no promotor
    \item Visualize a estrutura do gene (exons/introns)
\end{enumerate}
\end{alertblock}

\begin{block}{Questão 5: RNA-seq}
Com dados de contagens de RNA-seq:
\begin{enumerate}
    \item Normalize os dados (TPM ou DESeq2)
    \item Identifique genes diferencialmente expressos
    \item Faça análise de enriquecimento GO
    \item Crie volcano plot e heatmap
\end{enumerate}
\end{block}
\end{frame}

\begin{frame}{Exercícios - Parte 3}
\begin{block}{Questão 6: Predição de miRNA}
Dada uma sequência 3' UTR de um mRNA:
\begin{enumerate}
    \item Use TargetScan para predizer miRNAs que a regulam
    \item Analise a conservação evolutiva dos sítios
    \item Verifique dados de expressão: correlação negativa?
    \item Proponha experimento para validar a interação
\end{enumerate}
\end{block}

\begin{alertblock}{Projeto Final}
Desenvolva um pipeline de análise completo:
\begin{enumerate}
    \item Download de dados de RNA-seq do GEO
    \item Controle de qualidade (FastQC)
    \item Alinhamento ao genoma de referência
    \item Contagem de reads por gene
    \item Análise de expressão diferencial
    \item Análise de enriquecimento funcional
    \item Visualização e relatório
\end{enumerate}
\end{alertblock}
\end{frame}

% ============================================================================
% SEÇÃO 11: REFERÊNCIAS
% ============================================================================
\section{Referências}

\begin{frame}{Referências Principais}
\begin{thebibliography}{99}
\bibitem{watson} Watson J.D. et al. (2013). \textit{Molecular Biology of the Gene}, 7th ed. Pearson.

\bibitem{alberts} Alberts B. et al. (2022). \textit{Molecular Biology of the Cell}, 7th ed. Garland Science.

\bibitem{cock} Cock P.J.A. et al. (2009). Biopython: freely available Python tools for computational molecular biology and bioinformatics. \textit{Bioinformatics}.

\bibitem{love} Love M.I. et al. (2014). Moderated estimation of fold change and dispersion for RNA-seq data with DESeq2. \textit{Genome Biology}.

\bibitem{bartel} Bartel D.P. (2018). Metazoan MicroRNAs. \textit{Cell}.
\end{thebibliography}
\end{frame}

\begin{frame}{Leitura Complementar}
\begin{block}{Livros}
\begin{itemize}
    \item Lewin B. (2020). \textit{Genes XII}. Jones \& Bartlett.
    \item Lodish H. et al. (2021). \textit{Molecular Cell Biology}, 9th ed.
    \item Mount D.W. (2004). \textit{Bioinformatics: Sequence and Genome Analysis}.
\end{itemize}
\end{block}

\begin{block}{Recursos Online}
\begin{itemize}
    \item \textbf{NCBI Education:} \url{https://www.ncbi.nlm.nih.gov/guide/}
    \item \textbf{Ensembl:} \url{https://www.ensembl.org}
    \item \textbf{UCSC Genome Browser:} \url{https://genome.ucsc.edu}
    \item \textbf{GTEx Portal:} \url{https://gtexportal.org}
    \item \textbf{BioPython Tutorial:} \url{http://biopython.org/DIST/docs/tutorial/Tutorial.html}
\end{itemize}
\end{block}
\end{frame}

% ============================================================================
% SLIDE FINAL
% ============================================================================
\begin{frame}
\begin{center}
{\Huge Obrigado!}

\vspace{1cm}

{\Large Capítulo 6: Sistemas Gênicos}

\vspace{1cm}

\textbf{Próximo Capítulo:}\\
Protein Systems (Sistemas Proteicos)

\vspace{1cm}

\textit{Dúvidas? Perguntas?}
\end{center}
\end{frame}

\end{document}
